\documentclass[a4paper,12pt]{article}
\usepackage[utf8]{inputenc}
\usepackage[T1]{fontenc}
\usepackage[ngerman]{babel}
\usepackage{amsmath, amssymb}
\usepackage{geometry}
\usepackage{tikz}
\usetikzlibrary{shapes.geometric, arrows.meta, positioning, fit, backgrounds}
\usepackage{parskip}
\usepackage{titlesec}
\usepackage{booktabs}
\usepackage{enumitem}

\geometry{left=3cm, right=3cm, top=2.5cm, bottom=2.5cm}

\titleformat{\section}{\large\bfseries}{}{0em}{}[\titlerule]
\titleformat{\subsection}{\normalsize\bfseries}{}{0em}{}

\tikzset{
  box/.style={rectangle, draw, rounded corners=3pt, text width=4.2cm,
              align=center, minimum height=0.9cm, font=\small, inner sep=4pt},
  boxw/.style={box, text width=5.2cm},
  boxn/.style={box, text width=3.8cm},
  arrow/.style={-{Latex[length=2mm]}, thick},
}

\begin{document}

\begin{center}
  {\LARGE\bfseries Studierhinweise zu Lektion 2}\\[0.5em]
  {\large Analysis (Modul 61211)}
\end{center}

\vspace{1em}

Diese Kurseinheit dient der Übertragung von Eigenschaften stetiger reeller Funktionen
einer Veränderlichen auf Funktionen, die eine Teilmenge $M$ eines normierten Raumes
$(X, \|\cdot\|_X)$ in einen normierten Raum $(Y, \|\cdot\|_Y)$ abbilden. Dabei ist
$X = \mathbb{R}^n$ und $Y = \mathbb{R}^m$ der wichtigste Fall. Aus MG lässt sich
vieles fast wörtlich übernehmen und wird Ihnen vertraut vorkommen und rasch einleuchten.

Der zweite Abschnitt, in dem vor allem versucht wird, Möglichkeiten einer Veranschaulichung
von Funktionen zweier Veränderlicher zu skizzieren, bringt außer der Definition einiger
einfacher Funktionenklassen nicht viel Neues. Die Stetigkeit einer Funktion in einem
Punkt (Abschnitt~2.3) beruht auf dem Umgebungsbegriff, kann aber auch mithilfe von
Folgenkonvergenz ausgedrückt werden.

Die Abschnitte~2.4 und~2.5 behandeln Funktionen, die auf ihrem ganzen Definitionsbereich
$M$ stetig sind. Wichtiges Hilfsmittel ist die Charakterisierung dieser (globalen)
Stetigkeit mithilfe von offenen Mengen. Zentral in Abschnitt~2.4 ist der
"`Vererbungssatz"', wonach zusammenhängende Mengen von stetigen Funktionen in
zusammenhängende Mengen abgebildet werden. Im letzten Abschnitt geht es um die Frage,
wann sich Stetigkeit der Glieder einer Funktionenfolge auf die Grenzfunktion vererbt;
der wichtigste Satz ist der Satz von Weierstraß~2.6.5 über die Stetigkeit der
Grenzfunktion.

% -----------------------------------------------------------------------
\section*{Struktur der Kurseinheit 2}

\begin{center}
\begin{tikzpicture}[node distance=0.5cm and 0.6cm]

  % Column headers
  \node[font=\bfseries\small] (hA) at (0,0)    {Rückblick auf MG};
  \node[font=\bfseries\small] (hB) at (9.5,0)  {Ergänzungen und neue Inhalte};

  % ---- Left column: MG-Rückblick (Abschnitt 2.1) ----
  \node[box, below=0.8cm of hA] (k1) {2.1 Stetigkeit\\in einem Punkt};
  \node[box, below=0.4cm of k1] (k2) {2.1 Stetigkeitskriterien,\\Operationen};
  \node[box, below=0.4cm of k2] (k3) {2.1 stetiges Bild eines\\Intervalls, Zwischenwertsatz};
  \node[box, below=0.4cm of k3] (k4) {2.1 stetiges Bild von $[a,b]$};
  \node[box, below=0.4cm of k4] (k5) {2.1 Satz vom Maximum};
  \node[box, below=0.4cm of k5] (k6) {2.1 kompakte Menge in $\mathbb{R}$};
  \node[box, below=0.4cm of k6] (k7) {2.1 Überdeckungssatz\\von Heine--Borel};
  \node[box, below=0.4cm of k7] (k8) {2.1 gleichmäßige\\Stetigkeit};

  % ---- Right column: new content (2.2 - 2.6) ----
  \node[box, below=0.8cm of hB] (n1) {2.2 $f \in \mathrm{Abb}(M,Y)$,\\$M \subseteq X$; normierte Räume};
  \node[box, below=0.4cm of n1] (n2) {2.2 Veranschaulichung\\im Fall $X=\mathbb{R}^n$, $Y=\mathbb{R}^m$};
  \node[box, below=0.4cm of n2] (n3) {2.3 Stetigkeit\\in einem Punkt};
  \node[box, below=0.4cm of n3] (n4) {2.3 Stetigkeitskriterien,\\Operationen};
  \node[box, below=0.4cm of n4] (n5) {2.4 zusammenhängende\\Menge, stetiges Bild};
  \node[box, below=0.4cm of n5] (n6) {2.5 offene Überdeckung,\\kompakte Mengen};
  \node[box, below=0.4cm of n6] (n7) {2.5 stetiges Bild kompakter\\Mengen, Satz vom Maximum};
  \node[box, below=0.4cm of n7] (n8) {2.5 Heine--Borel in $\mathbb{R}^n$,\\glm.\ Stetigkeit};
  \node[box, below=0.4cm of n8] (n9) {2.6 punktweise und glm.\\Konvergenz, Funktionenreihen};
  \node[box, below=0.4cm of n9] (n10){2.6 Stetigkeit\\der Grenzfunktion};

  % Arrows left column
  \draw[arrow] (k1) -- (k2);
  \draw[arrow] (k2) -- (k3);
  \draw[arrow] (k3) -- (k4);
  \draw[arrow] (k4) -- (k5);
  \draw[arrow] (k5) -- (k6);
  \draw[arrow] (k6) -- (k7);
  \draw[arrow] (k7) -- (k8);

  % Arrows right column
  \draw[arrow] (n1) -- (n2);
  \draw[arrow] (n2) -- (n3);
  \draw[arrow] (n3) -- (n4);
  \draw[arrow] (n4) -- (n5);
  \draw[arrow] (n5) -- (n6);
  \draw[arrow] (n6) -- (n7);
  \draw[arrow] (n7) -- (n8);
  \draw[arrow] (n8) -- (n9);
  \draw[arrow] (n9) -- (n10);

  % Cross arrows (left -> right): Verallgemeinerung
  \draw[arrow, dashed] (k1.east) -- ++(0.3,0) |- (n3.west);
  \draw[arrow, dashed] (k3.east) -- ++(0.3,0) |- (n5.west);
  \draw[arrow, dashed] (k6.east) -- ++(0.3,0) |- (n6.west);
  \draw[arrow, dashed] (k7.east) -- ++(0.3,0) |- (n8.west);
  \draw[arrow, dashed] (k8.east) -- ++(0.3,0) |- (n8.west);

\end{tikzpicture}
\end{center}

\newpage

% -----------------------------------------------------------------------
\section*{Zielelement 2.1, Teil 1 -- Rückblick und Ergänzungen: Stetigkeit, lokale Eigenschaften}

\subsection*{Lerninhalte}

\begin{center}
\begin{tikzpicture}[node distance=0.5cm and 1.0cm]
  \node[box] (fkt) {Funktion: Definitions-\\bereich, Zielbereich,\\Argument, Funktionswert,\\Bild, Urbild};
  \node[box, right=1.0cm of fkt] (inj) {injektive Funktion,\\Umkehrfunktion};
  \node[box, below=0.5cm of fkt] (mon) {monotone Funktion};
  \node[box, right=1.0cm of mon] (umk) {Umkehrfunktion streng\\monotoner Funktionen};
  \node[box, below=0.5cm of mon] (ste) {in einem Punkt\\stetige Funktion};
  \node[box, right=1.0cm of ste] (lok) {lokale Beschränktheit,\\lokale Trennung};
  \node[box, below=0.5cm of ste] (kri) {$\varepsilon$-$\delta$-Kriterium,\\Folgenkriterium};
  \node[box, right=1.0cm of kri] (ops) {Operationen mit\\stetigen Funktionen};
  \node[box, below=0.5cm of ops] (pol) {Polynomfunktionen,\\rationale Funktionen};
  \draw[arrow] (fkt) -- (mon);
  \draw[arrow] (inj.south) |- (umk.north);
  \draw[arrow] (mon) -- (ste);
  \draw[arrow] (ste) -- (kri);
  \draw[arrow] (kri) -- (ops);
  \draw[arrow] (ste.east) -- (lok.west);
  \draw[arrow] (ops) -- (pol);
\end{tikzpicture}
\end{center}

\subsection*{Lernziele}

Nach Durcharbeiten des ersten Teils dieses Abschnitts sollten Sie sich wichtige
Begriffe und Zusammenhänge, die Sie in MG kennen gelernt haben, wieder
vergegenwärtigt haben und darüber hinaus in der Lage sein,
\begin{itemize}
  \item die Stetigkeit einer Funktion $f : M \to \mathbb{R}$ mit $M \subseteq \mathbb{R}$
        mithilfe des Umgebungsbegriffs auszudrücken,
  \item darzulegen, dass Stetigkeit eine lokale Eigenschaft ist.
\end{itemize}

\subsection*{Selbstkontrollelement 2.1, Teil 1}

Ist $f : \mathbb{R} \to \mathbb{R}$ eine Funktion mit
\[
  |f(x)| \leq |x| \quad \text{für jedes } x \in \mathbb{R},
\]
so ist $f$ in $0$ stetig. Können Sie dies mithilfe der Definition 2.1.6, mithilfe
des $\varepsilon$-$\delta$-Kriteriums 2.1.8, mithilfe des Folgenkriteriums 2.1.9
begründen?

\newpage

% -----------------------------------------------------------------------
\section*{Zielelement 2.1, Teil 2 -- Rückblick und Ergänzungen: Stetigkeit, globale Eigenschaften}

\subsection*{Lerninhalte}

\begin{center}
\begin{tikzpicture}[node distance=0.5cm and 1.0cm]
  \node[box] (intv) {auf einem Intervall\\stetige Funktion};
  \node[box, right=1.0cm of intv] (zws) {Zwischenwertsatz,\\stetiges Bild eines\\Intervalls};
  \node[box, below=0.5cm of intv] (hfp) {Häufungspunkt,\\abgeschlossene Menge};
  \node[box, right=1.0cm of hfp] (umk) {Stetigkeit der\\Umkehrfunktion};
  \node[box, below=0.5cm of hfp] (bw)  {Satz von\\Bolzano--Weierstraß};
  \node[box, below=0.5cm of bw] (kmp) {kompakte Teilmenge\\in $\mathbb{R}$};
  \node[box, right=1.0cm of kmp] (bild) {stetiges Bild einer\\kompakten Menge};
  \node[box, below=0.5cm of kmp] (hb) {Überdeckungssatz\\von Heine--Borel};
  \node[box, right=1.0cm of hb] (max) {Existenz von\\Minimum und Maximum};
  \node[box, below=0.5cm of hb] (gst) {gleichmäßige Stetigkeit\\auf einer kompakten Menge};
  \draw[arrow] (intv) -- (hfp);
  \draw[arrow] (intv.east) -- (zws.west);
  \draw[arrow] (zws.south) |- (umk.north);
  \draw[arrow] (hfp) -- (bw);
  \draw[arrow] (bw) -- (kmp);
  \draw[arrow] (kmp) -- (hb);
  \draw[arrow] (kmp.east) -- (bild.west);
  \draw[arrow] (bild.south) |- (max.north);
  \draw[arrow] (hb) -- (gst);
\end{tikzpicture}
\end{center}

\subsection*{Lernziele}

Nach Durcharbeiten des zweiten Teils dieses Abschnitts sollten Sie sich weitere
wichtige Begriffe und Zusammenhänge, die Sie in MG kennen gelernt haben, wieder
vergegenwärtigt haben und darüber hinaus in der Lage sein,
\begin{itemize}
  \item den Begriff des Häufungspunktes mittels Umgebungen auszudrücken und den
        Satz über die Existenz von Häufungspunkten zu formulieren,
  \item zu erklären, was eine kompakte Teilmenge von $\mathbb{R}$ ist,
  \item zu beschreiben, welche Eigenschaften das Bild einer kompakten Menge unter
        einer stetigen Funktion hat,
  \item den Begriff der gleichmäßigen Stetigkeit mit dem der (einfachen) Stetigkeit
        zu vergleichen und in diesem Zusammenhang die Bedeutung eines kompakten
        Definitionsbereiches aufzuzeigen.
\end{itemize}

\subsection*{Selbstkontrollelement 2.1, Teil 2}

Die Funktion $f : \,]0,1] \to \mathbb{R}$, $x \mapsto f(x) := \tfrac{1}{x}$ nimmt zwar
ihr Infimum an, aber nicht ihr Supremum. Ist also der Satz von Weierstraß~2.1.23
falsch?

\newpage

% -----------------------------------------------------------------------
\section*{Zielelement 2.2 -- Allgemeines über Funktionen von $\mathbb{R}^n$ nach $\mathbb{R}^m$}

\subsection*{Lerninhalte}

\begin{center}
\begin{tikzpicture}[node distance=0.5cm and 1.0cm]
  \node[box] (abb) {Funktionen $f : M \to \mathbb{R}^m$\\mit $M \subseteq \mathbb{R}^n$};
  \node[box, below=0.5cm of abb] (graph) {Graph einer Funktion,\\Veranschaulichungen\\(Höhenlinien, Graph in $\mathbb{R}^3$)};
  \node[box, below=0.5cm of graph] (opr) {Operationen in $\mathrm{Abb}(M, \mathbb{R}^m)$:\\Addition, skalare Multiplikation,\\(bei $m=1$) Multiplikation};
  \node[box, below=0.5cm of opr] (lin) {lineare Funktionen\\$f : \mathbb{R}^n \to \mathbb{R}^m$};
  \node[box, right=1.0cm of lin] (poly) {Polynomfunktionen,\\rationale Funktionen\\mehrerer Veränderlicher};
  \draw[arrow] (abb) -- (graph);
  \draw[arrow] (graph) -- (opr);
  \draw[arrow] (opr) -- (lin);
  \draw[arrow] (lin.east) -- (poly.west);
\end{tikzpicture}
\end{center}

\subsection*{Lernziele}

Nach Durcharbeiten dieses Abschnitts sollten Sie in der Lage sein,
\begin{itemize}
  \item allgemein etwas über Funktionen von $\mathbb{R}^n$ nach $\mathbb{R}^m$ zu
        sagen (Darstellungsmöglichkeiten in einfachen Fällen, Operieren mit solchen
        Funktionen),
  \item lineare Funktionen (Homomorphismen) von $\mathbb{R}^n$ nach $\mathbb{R}^m$
        sowie Polynomfunktionen und rationale Funktionen mehrerer Veränderlicher zu
        charakterisieren.
\end{itemize}

\subsection*{Selbstkontrollelement 2.2}

Sei $f : \mathbb{R}^2 \to \mathbb{R}^2$ die "`Drehung um den Nullpunkt im
Gegenuhrzeigersinn um den Winkel $\tfrac{\pi}{2}$"', d.\,h.\ um $90^\circ$. Dann ist
$f$ linear. Welches sind die Vektoren $a_1$ und $a_2$ aus $\mathbb{R}^2$, von denen
in 2.2.4(ii) die Rede ist?

\newpage

% -----------------------------------------------------------------------
\section*{Zielelement 2.3 -- Stetigkeit. Lokale Eigenschaften}

\subsection*{Lerninhalte}

\begin{center}
\begin{tikzpicture}[node distance=0.5cm and 1.0cm]
  \node[box] (ste) {in einem Punkt stetige\\Funktionen\\$f : M \to Y$, $M \subseteq X$};
  \node[box, right=1.0cm of ste] (lok) {lokale Trennung,\\lokale Beschränktheit};
  \node[box, below=0.5cm of ste] (kri) {$\varepsilon$-$\delta$-Kriterium,\\Folgenkriterium};
  \node[box, right=1.0cm of kri] (res) {Stetigkeit der\\Restriktion, Verkettung};
  \node[box, below=0.5cm of kri] (opr) {Operationen mit\\stetigen Funktionen};
  \node[box, right=1.0cm of opr] (bsp) {Beispiele: konstante,\\lineare, Polynom-,\\rationale Funktionen};
  \node[box, below=0.5cm of opr] (proj) {Projektionen,\\komponentenweise\\Stetigkeit};
  \draw[arrow] (ste) -- (kri);
  \draw[arrow] (ste.east) -- (lok.west);
  \draw[arrow] (kri) -- (opr);
  \draw[arrow] (kri.east) -- (res.west);
  \draw[arrow] (opr) -- (proj);
  \draw[arrow] (opr.east) -- (bsp.west);
\end{tikzpicture}
\end{center}

\subsection*{Lernziele}

Nach Durcharbeiten dieses Abschnitts sollten Sie in der Lage sein,
\begin{itemize}
  \item die Stetigkeit einer Funktion in einem Punkt zu definieren,
  \item das $\varepsilon$-$\delta$-Kriterium und das Folgenkriterium für Stetigkeit
        zu formulieren,
  \item Operationen mit stetigen Funktionen anzugeben, die wieder zu stetigen
        Funktionen führen,
  \item Beispiele stetiger Funktionen anzugeben, insbesondere die linearen Funktionen
        von $\mathbb{R}^n$ nach $\mathbb{R}^m$, sowie Polynomfunktionen und rationale
        Funktionen auf $\mathbb{R}^n$ zu beschreiben,
  \item den lokalen Charakter des Stetigkeitsbegriffs zu erläutern.
\end{itemize}

\subsection*{Selbstkontrollelement 2.3}

Können Sie zwei Funktionen $f$ und $g$ mit gleichem Definitionsbereich angeben, die
beide in einem Punkt $a$ ihres gemeinsamen Definitionsbereichs unstetig sind, deren
Summe $f+g$ aber in diesem Punkt $a$ stetig ist?

\newpage

% -----------------------------------------------------------------------
\section*{Zielelement 2.4 -- Stetige Funktionen auf zusammenhängenden Mengen}

\subsection*{Lerninhalte}

\begin{center}
\begin{tikzpicture}[node distance=0.5cm and 1.0cm]
  \node[box] (off) {offene, abgeschlossene\\Mengen in $(X, \|\cdot\|)$};
  \node[box, right=1.0cm of off] (glst) {Charakterisierung der\\globalen Stetigkeit mittels\\offener Mengen};
  \node[box, below=0.5cm of off] (zus) {zusammenhängende\\Mengen in $(X, \|\cdot\|)$};
  \node[box, right=1.0cm of zus] (bild) {stetiges Bild\\zusammenhängender Mengen};
  \node[box, below=0.5cm of zus] (zusr) {zusammenhängende\\Mengen in $\mathbb{R}$\\(= Intervalle)};
  \node[box, right=1.0cm of zusr] (zwv) {(verallgemeinerter)\\Zwischenwertsatz};
  \draw[arrow] (off) -- (zus);
  \draw[arrow] (off.east) -- (glst.west);
  \draw[arrow] (zus) -- (zusr);
  \draw[arrow] (zus.east) -- (bild.west);
  \draw[arrow] (bild) -- (zwv);
  \draw[arrow] (zusr.east) -- (zwv.west);
\end{tikzpicture}
\end{center}

\subsection*{Lernziele}

Nach Durcharbeiten dieses Abschnitts sollten Sie in der Lage sein,
\begin{itemize}
  \item offene Mengen in einem normierten Vektorraum zu definieren,
  \item die (globale) Stetigkeit einer Funktion mithilfe offener Mengen zu
        charakterisieren,
  \item den Begriff der zusammenhängenden Menge zu beschreiben und den Satz über
        das stetige Bild zusammenhängender Mengen zu formulieren,
  \item die zusammenhängenden Mengen im Raum $\mathbb{R}$ zu charakterisieren und den
        bekannten Zwischenwertsatz als Spezialfall des vorgenannten Satzes über das
        stetige Bild zusammenhängender Mengen herauszuarbeiten.
\end{itemize}

\subsection*{Selbstkontrollelement 2.4}

"`Geraden"' in $\mathbb{R}^n$, also Mengen
$G_{a,b} := \{a + t(b-a) \in \mathbb{R}^n \mid t \in \mathbb{R}\}$ mit
$a, b \in \mathbb{R}^n$ und $a \neq b$, sind zusammenhängende Teilmengen von
$\mathbb{R}^n$. Warum?

\newpage

% -----------------------------------------------------------------------
\section*{Zielelement 2.5 -- Stetige Funktionen auf kompakten Mengen}

\subsection*{Lerninhalte}

\begin{center}
\begin{tikzpicture}[node distance=0.5cm and 1.0cm]
  \node[box] (ueb) {offene Überdeckung von $M$,\\endliche Teilüberdeckung};
  \node[box, below=0.5cm of ueb] (kmp) {kompakte Menge $M$};
  \node[box, right=1.0cm of kmp] (fkmp) {folgenkompakte Menge $M$\\(äquivalent zu kompakt)};
  \node[box, below=0.5cm of kmp] (abg) {Kompaktheit und\\Abgeschlossenheit};
  \node[box, below=0.5cm of abg] (bsr) {Kompaktheit und\\Beschränktheit};
  \node[box, below=0.5cm of bsr] (hb)  {kompakte Teilmengen\\in $\mathbb{R}^n$ (Heine--Borel)};
  \node[box, below=0.5cm of hb] (bild) {stetiges Bild\\kompakter Mengen};
  \node[box, right=1.0cm of bild] (max) {Satz vom Maximum};
  \node[box, below=0.5cm of bild] (gst) {gleichmäßige Stetigkeit\\auf kompakten Mengen};
  \draw[arrow] (ueb) -- (kmp);
  \draw[arrow] (kmp.east) -- (fkmp.west);
  \draw[arrow] (kmp) -- (abg);
  \draw[arrow] (abg) -- (bsr);
  \draw[arrow] (bsr) -- (hb);
  \draw[arrow] (hb) -- (bild);
  \draw[arrow] (bild.east) -- (max.west);
  \draw[arrow] (bild) -- (gst);
\end{tikzpicture}
\end{center}

\subsection*{Lernziele}

Nach Durcharbeiten dieses Abschnitts sollten Sie in der Lage sein,
\begin{itemize}
  \item die Begriffe "`kompakte Menge"' und "`folgenkompakte Menge"' zu definieren
        und miteinander in Beziehung zu setzen,
  \item die Kompaktheit mit der Abgeschlossenheit und der Beschränktheit in
        Beziehung zu setzen,
  \item die kompakten Teilmengen in $\mathbb{R}^n$ zu charakterisieren,
  \item zwei wichtige Sätze über Funktionen, die auf einer kompakten Menge stetig
        sind, zu formulieren und im Einzelnen zu erläutern,
  \item den Weierstraßschen Satz vom Maximum aus dem Satz über das stetige Bild
        kompakter Mengen zu folgern.
\end{itemize}

\subsection*{Selbstkontrollelement 2.5}

Seien $(X, \|\cdot\|)$ ein normierter Raum und $\emptyset \neq M \subseteq X$.
Machen Sie sich die Beziehungen zwischen den Aussagen "`$M$ ist kompakt"',
"`$M$ ist folgenkompakt"', "`$M$ ist abgeschlossen"', "`$M$ ist beschränkt"' klar.

\newpage

% -----------------------------------------------------------------------
\section*{Zielelement 2.6 -- Punktweise und gleichmäßige Konvergenz von Funktionenfolgen}

\subsection*{Lerninhalte}

\begin{center}
\begin{tikzpicture}[node distance=0.5cm and 1.0cm]
  \node[box] (pkt) {punktweise Konvergenz\\von Funktionenfolgen\\bzw.\ Funktionenreihen};
  \node[box, right=1.0cm of pkt] (glm) {gleichmäßige Konvergenz\\von Funktionenfolgen\\bzw.\ Funktionenreihen};
  \node[box, below=0.5cm of pkt] (maj) {Majorantenkriterium für\\gleichmäßige Konvergenz};
  \node[box, right=1.0cm of maj] (ste) {Stetigkeit der Grenzfunktion\\bzw.\ der Summenfunktion\\(Satz von Weierstraß)};
  \draw[arrow] (pkt) -- (glm);
  \draw[arrow] (pkt) -- (maj);
  \draw[arrow] (glm) -- (ste);
  \draw[arrow] (maj) -- (ste);
\end{tikzpicture}
\end{center}

\subsection*{Lernziele}

Nach Durcharbeiten dieses Abschnitts sollten Sie in der Lage sein,
\begin{itemize}
  \item die Begriffe "`punktweise Konvergenz"' und "`gleichmäßige Konvergenz"' zu
        definieren und miteinander in Beziehung zu setzen,
  \item das Majorantenkriterium für gleichmäßige Konvergenz zu formulieren und an
        Beispielen zu erläutern,
  \item den Weierstraßschen Satz von der Stetigkeit der Grenzfunktion einer
        gleichmäßig konvergenten Funktionenfolge zu formulieren und anzuwenden.
\end{itemize}

\subsection*{Selbstkontrollelement 2.6}

"`Sehen"' Sie, dass die Folge $(p_k)$ mit $p_k : [0,1] \to \mathbb{R}$,
$x \mapsto p_k(x) := x^k$ zwar punktweise, aber nicht gleichmäßig konvergiert?

\end{document}
